Ma thèse propose une étude sur l'approximation de rang faible des matrices et
des tenseurs. Plus précisément nous avons définis de nouveau algorithmes
stables pour la décomposition de rang faibles pour les tenseurs. Et des
méthodes de précisions mixtes pour accélérer les algorithmes en calculant en
faibles précisions tout en atteignant une haute précision d'approximation.

La gestion des données est souvent réalisée par des objets mathématiques tels
que les matrices et les tenseurs, qui sont la généralisation des matrices à
plus de deux dimensions. Certains domaines d'application nécessitent de stocker
trop d'éléments, créant des tenseurs trop grands ; ce problème est connu sous
le nom de \emph{curse of dimensionality}. Des méthodes mathématiques telles que
les approximations de rang faible ont été développées pour réduire la
dimensionnalité de ces objets malgré un coût très élevé en temps de calcul. De
plus, de nouvelles architectures informatiques telles que les GPU nous
permettent d'effectuer des calculs rapidement, notamment lors de calculs en
faible précision. Combiner ces nouvelles architectures avec l'approximation de
rang faible est une solution malgré la qualité des résultats altérée par la
faible précision. Cette thèse vise à proposer des algorithmes d'approximation
de rang faible stables en faible précision tout en conservant l'accélération
inhérente au calcul en faible précision, ce qui est réalisable grâce au calcul
en précision mixte.

Nous avons développé une méthode générale d'approximation de tenseurs en
précision mixte en calculant d'abord une approximation en faible précision et
en l'affinant itérativement avec une précision supérieure pour maintenir la
qualité du résultat. Sachant que cette accélération provient principalement des
architectures GPU, plus précisément d'unités de calcul spécialisées appelées
\emph{tensor cores}, nous avons développé une méthode générale d'approximation
matricielle pour les architectures GPU en précision mixte utilisant ces
\emph{tensor cores}. Notre méthode maintient la qualité du résultat, mais au
prix d'une approximation de dimension supérieur à celle des applications
standards. Pour compenser cet écart, des méthodes de recompression de dimension
existent pour différents formats de tenseurs. Notre contribution finale propose
un cadre pour l'analyse de stabilité des opérations de réseaux tenseurs
communs, qui nous guide également vers une méthode de recompression stable.